\section{Problemas resueltos de diseño de experimentos y ANOVA}

\subsection{Enunciados de los problemas}

\subsubsection*{1. Nivel Fundamental (Sencillos / Directos)}

\begin{problema}
	\label{prob:anova_un_factor_ml}
	Un científico de datos desea comparar el rendimiento general (medido a través de la métrica $F_1$-score porcentual) de cuatro arquitecturas de clasificación de aprendizaje automático: Bosques Aleatorios (\emph{Random Forest}, Tratamiento 1), Máquinas de Vectores de Soporte (\emph{SVM}, Tratamiento 2), \emph{XGBoost} (Tratamiento 3) y Redes Neuronales Profundas (\emph{DNN}, Tratamiento 4).
	
	Para garantizar objetividad, entrena y evalúa los cuatro algoritmos de manera completamente aleatoria sobre 5 conjuntos de datos de referencia independientes y de alta dimensión ($n = 5$ réplicas por algoritmo, para un total de $N = 20$ experimentos). Las medias muestrales y desviaciones estándares de los $F_1$-scores porcentuales obtenidos en las 5 pruebas de cada modelo son:
	\begin{itemize}
		\item \textbf{Random Forest ($i=1$):} $\bar{Y}_{1.} = 84.2\%$, $S_1 = 3.1\%$.
		\item \textbf{SVM ($i=2$):} $\bar{Y}_{2.} = 79.8\%$, $S_2 = 3.4\%$.
		\item \textbf{XGBoost ($i=3$):} $\bar{Y}_{3.} = 88.6\%$, $S_3 = 2.8\%$.
		\item \textbf{Redes Neuronales ($i=4$):} $\bar{Y}_{4.} = 86.4\%$, $S_4 = 3.5\%$.
	\end{itemize}
	Asumiendo normalidad y homoscedasticidad, construya la tabla completa de Análisis de Varianza de un factor al nivel de significancia $\alpha = 0.05$. ¿Existe evidencia estadística para afirmar que al menos uno de los algoritmos difiere en su $F_1$-score promedio?
\end{problema}

\begin{problema}
	\label{prob:posthoc_tukey_ml}
	Dado que el ANOVA global del Problema \ref{prob:anova_un_factor_ml} resultó altamente significativo, aplique el \textbf{procedimiento post-hoc HSD de Tukey} con una significancia global por familia de $\alpha = 0.05$ para comparar simultáneamente las seis parejas posibles de algoritmos e identificar exactamente cuáles modelos difieren entre sí en su rendimiento medio.
\end{problema}

\begin{problema}
	\label{prob:anova_fund_particion}
	En una auditoría de control de calidad sobre la resistencia a la tracción (en MegaPascal, MPa) de un polímero sintético sintetizado en cuatro formulaciones químicas ($k = 4$), un ingeniero recopila una muestra aleatoria balanceada con una suma total de observaciones $N = 30$. Tras calcular la variabilidad en los datos, registra una Suma de Cuadrados Total $\text{SCT} = 520.0$ $\text{MPa}^2$ y una Suma de Cuadrados del Error Residual $\text{SCE} = 180.0$ $\text{MPa}^2$.
	\begin{enumerate}
		\item Calcule la Suma de Cuadrados de los Tratamientos ($\text{SCTR}$) aplicando el principio de aditividad del ANOVA de un factor.
		\item Determine los grados de libertad asociados al tratamiento ($\nu_{\text{TR}}$), al error residual ($\nu_E$) y al total ($\nu_T$).
		\item Calcule los Cuadrados Medios del Tratamiento ($\text{CMTR}$) y del Error ($\text{CME}$), así como el estadístico de prueba $F$.
		\item Verifique al nivel del 1\% ($\alpha = 0.01$) si las cuatro formulaciones producen una resistencia a la tracción estadísticamente diferente (el cuantil crítico en tablas es $F_{0.01, \, 3, \, 26} = 4.637$).
	\end{enumerate}
\end{problema}

\subsubsection*{2. Nivel Operativo (Intermedios / De Desarrollo)}

\begin{problema}
	\label{prob:dbca_servidores}
	Para evaluar el tiempo de procesamiento de consultas complejas SQL en segundos (métrica donde menor es mejor) en $k = 3$ motores de optimización en bases de datos relacionales (Motor A, Motor B y Motor C), un ingeniero advierte que el hardware de los servidores de la empresa es muy heterogéneo y no puede mezclarse sin control. 
	
	Por tanto, diseña un \textbf{Diseño en Bloques Completos al Azar (DBCA)} utilizando como factor de bloqueo las generaciones físicas del hardware de $b = 4$ servidores corporativos distintos (Servidor 1: Última Gen, Servidor 2: Gen Media, Servidor 3: Gen Antigua, Servidor 4: Gen Legacy virtualizada). Cada motor es evaluado exactamente una vez dentro de cada servidor, obteniéndose los datos de la Tabla \ref{tab:dbca_datos}.
	\begin{table}[h!]
		\centering
		\begin{tabular}{lccccc}
			\hline
			\textbf{Motor (Tratamiento $i$)} & \textbf{Serv. 1 ($j=1$)} & \textbf{Serv. 2 ($j=2$)} & \textbf{Serv. 3 ($j=3$)} & \textbf{Serv. 4 ($j=4$)} & \textbf{Media Trat. ($\bar{Y}_{i.}$)} \\ \hline
			Motor A ($i=1$) & 12.0 & 15.5 & 22.0 & 28.5 & \textbf{19.50} \\
			Motor B ($i=2$) & 8.5 & 11.0 & 16.5 & 22.0 & \textbf{14.50} \\
			Motor C ($i=3$) & 9.5 & 12.5 & 18.5 & 23.5 & \textbf{16.00} \\ \hline
			\textbf{Media Bloque ($\bar{Y}_{.j}$)} & \textbf{10.00} & \textbf{13.00} & \textbf{19.00} & \textbf{24.67} & $\bar{Y}_{..} = \textbf{16.67}$ \\ \hline
		\end{tabular}
		\caption{Tiempos de ejecución en segundos bajo un DBCA con 3 motores y 4 servidores.}
		\label{tab:dbca_datos}
	\end{table}
	
	\begin{enumerate}
		\item Construya la tabla completa de ANOVA para el Diseño en Bloques Completos al Azar con $\alpha = 0.05$.
		\item ¿Existe una diferencia significativa entre los tres motores de optimización?
		\item ¿Fue útil y efectivo desde el punto de vista de eficiencia estadística haber incorporado a los servidores como un factor de bloqueo en el diseño?
	\end{enumerate}
\end{problema}

\begin{problema}
	\label{prob:bonferroni_lsd}
	Una firma de ingeniería en la nube compara la latencia media de transferencia de datos en milisegundos bajo $k = 5$ protocolos de enrutamiento distintos. El ANOVA arrojó un Cuadrado Medio del Error $\text{CME} = 18.0$ $\text{ms}^2$ con $\nu_E = 25$ grados de libertad en el residual, y cada protocolo fue evaluado con $n = 6$ réplicas.
	\begin{enumerate}
		\item Calcule la Diferencia Mínima Significativa (\textbf{LSD de Fisher}) al nivel $\alpha = 0.05$.
		\item Calcule el margen crítico de diferencia por el método de la \textbf{Corrección de Bonferroni} para mantener una tasa de error global máxima de $\alpha = 0.05$.
		\item Compare cuantitativamente ambos umbrales y explique qué implicaciones tiene la diferencia en la sensibilidad de detección en un proyecto de ciencia de datos.
	\end{enumerate}
\end{problema}

\begin{problema}
	\label{prob:anova_oper_desbalanceado_r2}
	En una investigación sobre eficiencia computacional, se mide el consumo del ancho de banda de CPU en porcentaje ($\%$) al ejecutar procesos intensivos en tres entornos de virtualización diferentes ($k = 3$). Debido a interrupciones en la granja de servidores, las muestras recopiladas quedaron desbalanceadas:
	\begin{itemize}
		\item \textbf{Entorno 1:} $n_1 = 6$ réplicas, media $\bar{Y}_{1.} = 45.0\%$, varianza muestral $S_1^2 = 12.0$.
		\item \textbf{Entorno 2:} $n_2 = 8$ réplicas, media $\bar{Y}_{2.} = 52.0\%$, varianza muestral $S_2^2 = 15.0$.
		\item \textbf{Entorno 3:} $n_3 = 5$ réplicas, media $\bar{Y}_{3.} = 38.0\%$, varianza muestral $S_3^2 = 10.0$.
	\end{itemize}
	\begin{enumerate}
		\item Calcule la gran media general $\bar{Y}_{..}$ ponderando correctamente los tamaños muestrales ($N = 19$).
		\item Calcule la Suma de Cuadrados de los Tratamientos ($\text{SCTR}$), la Suma de Cuadrados del Error ($\text{SCE}$) y construya la tabla ANOVA para comprobar si el consumo medio de CPU difiere al nivel $\alpha = 0.05$ ($F_{0.05, \, 2, \, 16} = 3.634$).
		\item Calcule el coeficiente de determinación múltiple del ANOVA ($R^2 = \dfrac{\text{SCTR}}{\text{SCT}}$) y el coeficiente $R^2$ ajustado, interpretando qué porcentaje de la varianza en el consumo de CPU es explicado por el tipo de virtualización.
	\end{enumerate}
\end{problema}

\subsubsection*{3. Nivel Analítico (Avanzados / De Conexión)}

\begin{problema}
	\label{prob:verificacion_supuestos_anova}
	Al auditar un ANOVA de un factor con $k = 3$ configuraciones de memoria caché y tamaños muestrales pequeños ($n_1=6, n_2=6, n_3=6$), las varianzas muestrales calculadas son $S_1^2 = 1.2$, $S_2^2 = 1.5$ y $S_3^2 = 8.9$.
	
	Un analista joven propone proceder con el ANOVA estándar aduciendo que "el ANOVA es robusto a violaciones en las varianzas si los tamaños grupales son iguales". ¿Es correcta esta afirmación ante estos datos? Explique qué procedimiento de verificación debe realizarse cuantitativamente y qué alternativa metodológica procedería si el supuesto falla.
\end{problema}

\begin{problema}
	\label{prob:potencia_tamano_anova}
	Un equipo de ciencia de datos planea un nuevo experimento (ANOVA de un factor) para comparar la velocidad de ejecución en peticiones de API en $k = 4$ lenguajes de backend nativos. Por estudios históricos del sistema, se estima que la varianza residual intrínseca de las peticiones es de $\sigma^2 = 16$ $\text{ms}^2$ ($\sigma = 4$ ms).
	
	El equipo desea determinar el tamaño de muestra de réplicas $n$ por lenguaje para que el ANOVA tenga una \textbf{potencia estadística del $80\%$ ($1 - \beta = 0.80$)} para detectar al nivel nominal de $\alpha = 0.05$ una diferencia máxima de al menos $\Delta = \mu_{\text{máx}} - \mu_{\text{mín}} = 5$ ms entre el lenguaje más rápido y el más lento del estudio. Utilice la aproximación de la variabilidad máxima mediante el parámetro $\Phi$ de las tablas de potencia de Pearson-Hartley o la fórmula del efecto estandarizado de Cohen $f$.
\end{problema}

\subsubsection*{4. Nivel Desafiante (Expertos / Deducción Profunda)}

\begin{problema}
	\label{prob:anova_desaf_cochran}
	\textbf{Deducción analítica de la partición cuadrática e independencia de formas cuadráticas por el Teorema de Cochran:}
	Sea un experimento unifactorial completamente aleatorizado de $k$ tratamientos y $n$ réplicas balanceadas ($N = nk$), bajo el modelo lineal aditivo normal:
	\begin{align}
		Y_{ij} = \mu + \tau_i + \varepsilon_{ij}, \quad i=1, \dots, k, \; j=1, \dots, n,
	\end{align}
	donde $\sum_{i=1}^k \tau_i = 0$ y los errores son i.i.d. $\varepsilon_{ij} \sim N(0, \sigma^2)$.
	\begin{enumerate}
		\item Partiendo de la identidad elemental de descomposición geométrica del error observacional:
		\begin{align}
			Y_{ij} - \bar{Y}_{..} = (\bar{Y}_{i.} - \bar{Y}_{..}) + (Y_{ij} - \bar{Y}_{i.}),
		\end{align}
		demuestre algebraicamente que al elevar ambos miembros al cuadrado y sumar sobre todos los $i$ y $j$, el doble producto cruzado se anula en forma exacta ($\sum_{i=1}^k \sum_{j=1}^n (\bar{Y}_{i.} - \bar{Y}_{..})(Y_{ij} - \bar{Y}_{i.}) \equiv 0$), probando el teorema de Pitágoras estadístico: $\text{SCT} = \text{SCTR} + \text{SCE}$.
		\item Demuestre que bajo la hipótesis nula universal $H_0: \tau_i = 0 \; \forall i$, las variables aleatorias normalizadas $\dfrac{\text{SCTR}}{\sigma^2}$ y $\dfrac{\text{SCE}}{\sigma^2}$ siguen distribuciones Chi-cuadrada independientes con $k-1$ y $N-k$ grados de libertad, respectivamente.
		\item Invoque el \textbf{Teorema de William G. Cochran} sobre la descomposición de la matriz identidad en matrices idempotentes ortogonales para justificar por qué los Cuadrados Medios $\text{CMTR}$ y $\text{CME}$ son independientes en probabilidad, concluyendo rigurosamente que el cociente $F = \dfrac{\text{CMTR}}{\text{CME}}$ sigue una distribución $F$ de Snedecor central de orden $(k-1, N-k)$.
	\end{enumerate}
\end{problema}

\begin{problema}
	\label{prob:anova_desaf_esperanzas_dbca}
	\textbf{Deducción del Valor Esperado de los Cuadrados Medios ($E[\text{CM}]$) en el DBCA y demostración matemática de la Eficiencia Relativa del bloqueo:}
	Considere un Diseño en Bloques Completos al Azar sin interacción con $k$ tratamientos y $b$ bloques ($N = kb$), modelado por:
	\begin{align}
		Y_{ij} = \mu + \tau_i + \beta_j + \varepsilon_{ij}, \quad \varepsilon_{ij} \sim N(0, \sigma^2 \mathbf{I}), \quad \sum_{i=1}^k \tau_i = 0, \quad \sum_{j=1}^b \beta_j = 0.
	\end{align}
	\begin{enumerate}
		\item Utilizando el operador esperanza matemática ($E[\cdot]$) y las propiedades de independencia y varianza de los errores $\varepsilon_{ij}$, demuestre paso a paso que las esperanzas del Cuadrado Medio de Tratamiento ($\text{CMTR}$), del Cuadrado Medio de Bloques ($\text{CMB}$) y del Cuadrado Medio del Error ($\text{CME}$) cumplen las identidades paramétricas exactas:
		\begin{align}
			E[\text{CMTR}] &= \sigma^2 + \frac{b}{k-1}\sum_{i=1}^k \tau_i^2, \\
			E[\text{CMB}] &= \sigma^2 + \frac{k}{b-1}\sum_{j=1}^b \beta_j^2, \\
			E[\text{CME}] &= \sigma^2.
		\end{align}
		\item En un Diseño Completamente Aleatorizado (DCA) que no reconoce los bloques pero utiliza las mismas $N = kb$ unidades experimentales, la varianza residual estimada se denota por $S_{\text{DCA}}^2$. Deduzca a partir de las esperanzas anteriores la fórmula de \textbf{Eficiencia Relativa ($\text{ER}$)} del DBCA respecto al DCA:
		\begin{align}
			S_{\text{DCA}}^2 \approx \frac{(b-1)\text{CMB} + b(k-1)\text{CME}}{kb - 1},
		\end{align}
		y demuestre algebraicamente cuál es la condición analítica sobre el Cuadrado Medio de los Bloques ($\text{CMB}$) que garantiza que el diseño bloqueado sea cuantitativamente superior en la precisión de los contrastes ($\text{ER} > 1$).
	\end{enumerate}
\end{problema}


\subsection{Soluciones detalladas}

\subsubsection*{1. Nivel Fundamental (Sencillos / Directos)}

\begin{solucion}
	\textbf{Paso 1: Planteamiento formal de hipótesis.}
	\begin{align}
		H_0&: \mu_1 = \mu_2 = \mu_3 = \mu_4 \quad (\text{Los cuatro algoritmos tienen idéntico rendimiento medio}). \\
		H_a&: \mu_i \neq \mu_j \text{ para al menos un par } (i, j).
	\end{align}
	
	\textbf{Paso 2: Cálculo de la gran media general ($\bar{Y}_{..}$).}
	Como los tamaños muestrales son balanceados ($n_1 = n_2 = n_3 = n_4 = 5$), la media general es el promedio aritmético simple de las 4 medias de tratamiento:
	\begin{align}
		\bar{Y}_{..} = \frac{84.2 + 79.8 + 88.6 + 86.4}{4} = \frac{339.0}{4} = 84.75\%.
	\end{align}
	
	\textbf{Paso 3: Cálculo de la Suma de Cuadrados del Tratamiento ($\text{SCTR}$).}
	La fórmula para $k = 4$ tratamientos es $\text{SCTR} = \sum_{i=1}^4 n (\bar{Y}_{i.} - \bar{Y}_{..})^2$:
	\begin{align}
		\text{SCTR} &= 5 \left[ (84.2 - 84.75)^2 + (79.8 - 84.75)^2 + (88.6 - 84.75)^2 + (86.4 - 84.75)^2 \right] \\
		&= 5 \left[ (-0.55)^2 + (-4.95)^2 + (3.85)^2 + (1.65)^2 \right] \\
		&= 5 \left[ 0.3025 + 24.5025 + 14.8225 + 2.7225 \right] = 5 \left[ 42.35 \right] = 211.75 \text{ (\%}^2\text{)}.
	\end{align}
	
	\textbf{Paso 4: Cálculo de la Suma de Cuadrados del Error ($\text{SCE}$).}
	Utilizamos la relación $\text{SCE} = \sum_{i=1}^4 (n_i - 1)S_i^2$. Como $n_i - 1 = 5 - 1 = 4$ en los cuatro casos:
	\begin{align}
		\text{SCE} &= 4 \left[ S_1^2 + S_2^2 + S_3^2 + S_4^2 \right] = 4 \left[ (3.1)^2 + (3.4)^2 + (2.8)^2 + (3.5)^2 \right] \\
		&= 4 \left[ 9.61 + 11.56 + 7.84 + 12.25 \right] = 4 \left[ 41.26 \right] = 165.04 \text{ (\%}^2\text{)}.
	\end{align}
	
	La Suma de Cuadrados Total resulta: $\text{SCT} = \text{SCTR} + \text{SCE} = 211.75 + 165.04 = 376.79$.
	
	\textbf{Paso 5: Grados de libertad y Cuadrados Medios ($\text{CM}$).}
	\begin{itemize}
		\item Grados de libertad del tratamiento: $\nu_{\text{TR}} = k - 1 = 4 - 1 = 3$.
		\item Grados de libertad del error: $\nu_E = N - k = 20 - 4 = 16$.
		\item Grados de libertad totales: $\nu_T = N - 1 = 19$.
	\end{itemize}
	Calculamos los cuadrados medios dividiendo:
	\begin{align}
		\text{CMTR} &= \frac{\text{SCTR}}{k - 1} = \frac{211.75}{3} \approx 70.5833, \\
		\text{CME} &= \frac{\text{SCE}}{N - k} = \frac{165.04}{16} = 10.315.
	\end{align}
	
	\textbf{Paso 6: Cálculo del estadístico $F$ observacional.}
	\begin{align}
		F = \frac{\text{CMTR}}{\text{CME}} = \frac{70.5833}{10.315} \approx 6.8428.
	\end{align}
	
	Sintetizamos toda la información en la Tabla \ref{tab:anova_sol_1}.
	\begin{table}[h!]
		\centering
		\begin{tabular}{lcccc}
			\hline
			\textbf{Fuente de Variación} & \textbf{SC} & \textbf{g.l. ($\nu$)} & \textbf{CM} & \textbf{Estadístico $F$} \\ \hline
			Tratamientos (Algoritmos) & 211.75 & 3 & 70.583 & \textbf{6.843} \\
			Error residual & 165.04 & 16 & 10.315 & \\ \hline
			\textbf{Total} & 376.79 & 19 & & \\ \hline
		\end{tabular}
		\caption{Tabla de ANOVA para la comparación de los cuatro algoritmos de clasificación.}
		\label{tab:anova_sol_1}
	\end{table}
	
	\textbf{Paso 7: Regla de decisión y conclusión.}
	Para un nivel de significancia $\alpha = 0.05$, consultamos la tabla de la distribución de Fisher-Snedecor con $\nu_1 = 3$ grados de libertad en el numerador y $\nu_2 = 16$ en el denominador:
	\begin{align}
		F_{0.05, \, 3, \, 16} = 3.239.
	\end{align}
	Dado que nuestro estadístico calculado $F = 6.843 \gg 3.239$ ($p-\text{valor} \approx 0.0035$), \textbf{rechazamos categóricamente la hipótesis nula $H_0$}. Demostramos con más del $99.6\%$ de confianza que existe una diferencia altamente significativa en la precisión media ($F_1$-score) de al menos uno de los cuatro algoritmos de aprendizaje automático sobre los conjuntos de datos evaluados.
\end{solucion}

\begin{solucion}
	\textbf{Paso 1: Obtención del cuantil del rango estandarizado ($q$).}
	De la tabla ANOVA del Problema \ref{prob:anova_un_factor_ml}, tenemos que el Cuadrado Medio del Error es $\text{CME} = 10.315$ con $\nu_E = 16$ grados de libertad, y contamos con $k = 4$ tratamientos. Todos con igual número de réplicas $n = 5$.
	
	Consultando las tablas estadísticas del rango estandarizado estudiantizado $q_{\alpha, k, \nu_E}$ para $\alpha = 0.05$, $k = 4$ y $\nu_E = 16$:
	\begin{align}
		q_{0.05, \, 4, \, 16} = 4.046.
	\end{align}
	
	\textbf{Paso 2: Cálculo de la Diferencia Honestamente Significativa ($\text{HSD}$).}
	Sustituyendo en la fórmula fundamental de Tukey para muestras balanceadas:
	\begin{align}
		\text{HSD} = q_{\alpha, \, k, \, \nu_E} \sqrt{\frac{\text{CME}}{n}} = 4.046 \sqrt{\frac{10.315}{5}} = 4.046 \sqrt{2.063} = 4.046 \times 1.4363 \approx 5.811\%.
	\end{align}
	Por consiguiente, \textbf{cualquier par de algoritmos cuya media difiera en valor absoluto por más de $5.811\%$ se considerará estadísticamente diferente} con un control estricto del error familiar al $5\%$.
	
	\textbf{Paso 3: Matriz ordenada de diferencias de medias grupales.}
	Ordenemos las medias muestrales de menor a mayor para sistematizar las comparaciones:
	\begin{itemize}
		\item $\bar{Y}_{\text{SVM}} = 79.8\%$
		\item $\bar{Y}_{\text{Random Forest (RF)}} = 84.2\%$
		\item $\bar{Y}_{\text{Redes Neuronales (DNN)}} = 86.4\%$
		\item $\bar{Y}_{\text{XGBoost (XGB)}} = 88.6\%$
	\end{itemize}
	Evaluamos las 6 diferencias por pares contra el umbral crítico $\text{HSD} = 5.811\%$:
	\begin{enumerate}
		\item \textbf{XGBoost vs SVM:} $|\bar{Y}_{\text{XGB}} - \bar{Y}_{\text{SVM}}| = |88.6 - 79.8| = 8.80\% > 5.811\%$ $\implies$ \textbf{Significativo ($*$)}.
		\item \textbf{XGBoost vs Random Forest:} $|\bar{Y}_{\text{XGB}} - \bar{Y}_{\text{RF}}| = |88.6 - 84.2| = 4.40\% \leq 5.811\%$ $\implies$ No significativo.
		\item \textbf{XGBoost vs Redes Neuronales:} $|\bar{Y}_{\text{XGB}} - \bar{Y}_{\text{DNN}}| = |88.6 - 86.4| = 2.20\% \leq 5.811\%$ $\implies$ No significativo.
		\item \textbf{Redes Neuronales vs SVM:} $|\bar{Y}_{\text{DNN}} - \bar{Y}_{\text{SVM}}| = |86.4 - 79.8| = 6.60\% > 5.811\%$ $\implies$ \textbf{Significativo ($*$)}.
		\item \textbf{Redes Neuronales vs Random Forest:} $|\bar{Y}_{\text{DNN}} - \bar{Y}_{\text{RF}}| = |86.4 - 84.2| = 2.20\% \leq 5.811\%$ $\implies$ No significativo.
		\item \textbf{Random Forest vs SVM:} $|\bar{Y}_{\text{RF}} - \bar{Y}_{\text{SVM}}| = |84.2 - 79.8| = 4.40\% \leq 5.811\%$ $\implies$ No significativo.
	\end{enumerate}
	\textbf{Conclusión técnica del análisis post-hoc:} El procedimiento de Tukey HSD revela y aisla con precisión la causa de significancia del ANOVA: \textbf{tanto XGBoost ($88.6\%$) como las Redes Neuronales Profundas ($86.4\%$) superan de manera estadísticamente probada a SVM ($79.8\%$)} en el $F_1$-score. Entre XGBoost, Redes Neuronales y Bosques Aleatorios no existen diferencias estadísticamente discernibles para el tamaño de muestra del estudio.
\end{solucion}

\begin{solucion}
	\begin{enumerate}
		\item \textbf{Cálculo de la Suma de Cuadrados de Tratamientos ($\text{SCTR}$):}
		Por el teorema de aditividad cuadrática ($\text{SCT} = \text{SCTR} + \text{SCE}$), despejamos de forma elemental:
		\begin{align}
			\text{SCTR} = \text{SCT} - \text{SCE} = 520.0 - 180.0 = 340.0 \text{ MPa}^2.
		\end{align}
		
		\item \textbf{Determinación de los grados de libertad ($\nu$):}
		Dado que el experimento abarca $k = 4$ formulaciones y $N = 30$ observaciones totales:
		\begin{align}
			\nu_{\text{TR}} &= k - 1 = 4 - 1 = 3. \\
			\nu_E &= N - k = 30 - 4 = 26. \\
			\nu_T &= N - 1 = 30 - 1 = 29.
		\end{align}
		Verificamos la aditividad del espacio de dimensión: $3 + 26 = 29$.
		
		\item \textbf{Cálculo de los Cuadrados Medios ($\text{CM}$) y estadístico $F$:}
		\begin{align}
			\text{CMTR} &= \frac{\text{SCTR}}{\nu_{\text{TR}}} = \frac{340.0}{3} \approx 113.333 \text{ MPa}^2. \\
			\text{CME} &= \frac{\text{SCE}}{\nu_E} = \frac{180.0}{26} \approx 6.923 \text{ MPa}^2.
		\end{align}
		El cociente de varianzas intergrupal vs. intragrupal arroja el estadístico observacional:
		\begin{align}
			F = \frac{\text{CMTR}}{\text{CME}} = \frac{113.333}{6.923} \approx 16.370.
		\end{align}
		
		\item \textbf{Evaluación docimásica y conclusión:}
		El cuantil crítico en la distribución $F$ con $3$ y $26$ grados de libertad al $1\%$ es $F_{0.01, \, 3, \, 26} = 4.637$.
		Puesto que nuestro estadístico muestral excede por un margen abrumador a la frontera ($16.370 \gg 4.637$), \textbf{se rechaza la hipótesis nula de igualdad de resistencias al nivel de significancia del 1\%}. El tipo de formulación química altera decisiva y significativamente la resistencia mecánica del polímero.
	\end{enumerate}
\end{solucion}

\subsubsection*{2. Nivel Operativo (Intermedios / De Desarrollo)}

\begin{solucion}
	\textbf{Parte 1: Partición de Sumas de Cuadrados y Tabla ANOVA DBCA.}
	La media general o gran media es $\bar{Y}_{..} = 16.6667$ s ($N = k \times b = 3 \times 4 = 12$ observaciones).
	
	\textbf{Suma de Cuadrados del Tratamiento ($\text{SCTR}$, $k=3$ con $b=4$ réplicas por motor):}
	\begin{align}
		\text{SCTR} &= b \sum_{i=1}^3 (\bar{Y}_{i.} - \bar{Y}_{..})^2 = 4 \left[ (19.5 - 16.6667)^2 + (14.5 - 16.6667)^2 + (16.0 - 16.6667)^2 \right] \nonumber \\
		&= 4 \left[ (2.8333)^2 + (-2.1667)^2 + (-0.6667)^2 \right] = 4 \left[ 8.0278 + 4.6944 + 0.4444 \right] \nonumber \\
		&= 4 \times 13.1667 = 52.6667 \text{ s}^2.
	\end{align}
	
	\textbf{Suma de Cuadrados de los Bloques ($\text{SCB}$, $b=4$ con $k=3$ réplicas por servidor):}
	\begin{align}
		\text{SCB} &= k \sum_{j=1}^4 (\bar{Y}_{.j} - \bar{Y}_{..})^2 \nonumber \\
		&= 3 \left[ (10.0 - 16.6667)^2 + (13.0 - 16.6667)^2 + (19.0 - 16.6667)^2 + (24.6667 - 16.6667)^2 \right] \nonumber \\
		&= 3 \left[ (-6.6667)^2 + (-3.6667)^2 + (2.3333)^2 + (8.0000)^2 \right] \nonumber \\
		&= 3 \left[ 44.4444 + 13.4444 + 5.4444 + 64.0000 \right] = 3 \times 127.3333 = 382.0000 \text{ s}^2.
	\end{align}
	
	\textbf{Suma de Cuadrados Total ($\text{SCT}$):} Sumamos las desviaciones de las 12 celdas respecto a $\bar{Y}_{..} = 16.6667$:
	\begin{align}
		\text{SCT} &= \sum_{i=1}^3 \sum_{j=1}^4 (Y_{ij} - 16.6667)^2 = (12-16.6667)^2 + (15.5-16.6667)^2 + \ldots + (23.5-16.6667)^2 \nonumber \\
		&= 21.78 + 1.36 + 28.44 + 140.03 + 66.70 + 32.11 + 0.03 + 28.44 + 51.36 + 17.36 + 3.36 + 46.70 \nonumber \\
		&= 437.6667 \text{ s}^2.
	\end{align}
	
	\textbf{Suma de Cuadrados del Error Residual ($\text{SCE}$):}
	Por diferencia ortogonal exacta:
	\begin{align}
		\text{SCE} = \text{SCT} - \text{SCTR} - \text{SCB} = 437.6667 - 52.6667 - 382.0000 = 3.0000 \text{ s}^2.
	\end{align}
	
	\textbf{Grados de libertad y Cuadrados Medios:}
	\begin{itemize}
		\item Tratamientos: $\nu_{\text{TR}} = k - 1 = 3 - 1 = 2 \implies \text{CMTR} = \frac{52.6667}{2} = 26.3333$.
		\item Bloques: $\nu_{\text{B}} = b - 1 = 4 - 1 = 3 \implies \text{CMB} = \frac{382.0000}{3} = 127.3333$.
		\item Error residual: $\nu_E = (k-1)(b-1) = 2 \times 3 = 6 \implies \text{CME} = \frac{3.0000}{6} = 0.5000$.
	\end{itemize}
	
	\textbf{Estadísticos $F$:}
	\begin{align}
		F_{\text{TR}} &= \frac{\text{CMTR}}{\text{CME}} = \frac{26.3333}{0.5000} = 52.667, \\
		F_{\text{B}} &= \frac{\text{CMB}}{\text{CME}} = \frac{127.3333}{0.5000} = 254.667.
	\end{align}
	Sintetizamos en la Tabla \ref{tab:anova_dbca_sol}.
	\begin{table}[h!]
		\centering
		\begin{tabular}{lcccc}
			\hline
			\textbf{Fuente de Variación} & \textbf{SC} & \textbf{g.l. ($\nu$)} & \textbf{CM} & \textbf{Estadístico $F$} \\ \hline
			Tratamientos (Motores SQL) & 52.667 & 2 & 26.333 & \textbf{52.667} \\
			Bloques (Servidores Hardware) & 382.000 & 3 & 127.333 & \textbf{254.667} \\
			Error residual & 3.000 & 6 & 0.500 & \\ \hline
			\textbf{Total} & 437.667 & 11 & & \\ \hline
		\end{tabular}
		\caption{Tabla ANOVA para el experimento DBCA de motores SQL e infraestructura de servidores.}
		\label{tab:anova_dbca_sol}
	\end{table}
	
	\textbf{Parte 2: Prueba sobre los tratamientos ($H_{0,\text{TR}}: \tau_1 = \tau_2 = \tau_3 = 0$).}
	Para el nivel de $\alpha = 0.05$ y grados de libertad $\nu_1 = 2, \nu_2 = 6$, el cuantil de rechazo es:
	\begin{align}
		F_{0.05, \, 2, \, 6} = 5.143.
	\end{align}
	Como el estadístico de tratamientos $F_{\text{TR}} = 52.667 \gg 5.143$, \textbf{rechazamos $H_{0,\text{TR}}$}. Existe una diferencia extremadamente significativa en el tiempo de procesamiento entre los tres motores SQL (observándose con claridad en las medias que el \textbf{Motor B} con $\bar{Y}_{2.} = 14.50$ s es el más rápido de los tres).
	
	\textbf{Parte 3: Evaluación de la efectividad del bloqueo ($H_{0,\text{B}}: \beta_1 = \ldots = \beta_4 = 0$).}
	Para el factor de bloques, consultamos el cuantil de Fisher con $\nu_1 = 3, \nu_2 = 6$:
	\begin{align}
		F_{0.05, \, 3, \, 6} = 4.757.
	\end{align}
	Dado que $F_{\text{B}} = 254.667 \gg 4.757$, la prueba sobre los bloques es colosalmente significativa. 
	
	\textbf{Conclusión metodológica sobre el diseño:} Haber incorporado las generaciones de los servidores como bloques fue una decisión metodológica \textbf{extraordinariamente exitosa y brillante}. El bloqueo logró absorber y aislar $382.0$ unidades de varianza pura de infraestructura ($\text{SCB}$). Si no hubiéramos bloqueado el diseño y hubiéramos evaluado los datos como un ANOVA de un factor normal, esos $382$ puntos de varianza habrían terminado sumados directamente al error residual ($\text{SCE}_{\text{un factor}} = 382 + 3 = 385$), lo que habría arrojado una nueva varianza residual $\text{CME}_{\text{un factor}} = \frac{385}{12-3} = 42.77$. 
	
	Bajo ese fallido análisis de un factor sin bloquear, el estadístico de los motores se habría desplomado a un insignificante:
	\begin{align}
		F_{\text{un factor}} = \frac{26.333}{42.77} = 0.615 \quad (\text{que es mucho menor al crítico } F_{0.05, 2, 9} = 4.26).
	\end{align}
	El bloqueo \textbf{rescató por completo el experimento de un falso negativo masivo (error Tipo II)} al purgar el ruido del hardware.
\end{solucion}

\begin{solucion}
	\textbf{Parte 1: Cálculo del LSD de Fisher.}
	Para $k = 5$ tratamientos, el número total de comparaciones por pares es $m = \binom{5}{2} = \frac{5 \times 4}{2} = 10$.
	El cuantil de la distribución $t$ de Student para una prueba individual de dos colas con $\alpha = 0.05$ ($\alpha/2 = 0.025$) y $\nu_E = 25$ grados de libertad es:
	\begin{align}
		t_{0.025, \, 25} = 2.060.
	\end{align}
	Sustituyendo en la fórmula del LSD:
	\begin{align}
		\text{LSD} = t_{0.025, \, 25} \sqrt{\text{CME} \left( \frac{1}{n} + \frac{1}{n} \right)} &= 2.060 \sqrt{18.0 \left( \frac{2}{6} \right)} = 2.060 \sqrt{6.0} \\
		&= 2.060 \times 2.4495 \approx 5.046 \text{ ms}.
	\end{align}
	
	\textbf{Parte 2: Cálculo con la Corrección de Bonferroni.}
	Para controlar el error familiar en $10$ comparaciones al nivel de $\alpha = 0.05$, el nivel de significancia individual de cada prueba se ajusta rígidamente a:
	\begin{align}
		\alpha^* = \frac{\alpha}{m} = \frac{0.05}{10} = 0.005 \implies \frac{\alpha^*}{2} = 0.0025.
	\end{align}
	Consultando las tablas de la distribución $t$ con $\nu_E = 25$ para el cuantil que deja una cola superior de $0.0025$:
	\begin{align}
		t_{0.0025, \, 25} = 3.078.
	\end{align}
	La diferencia crítica de Bonferroni resulta en:
	\begin{align}
		\text{CD}_{\text{Bonferroni}} = t_{0.0025, \, 25} \sqrt{\frac{2 \text{CME}}{n}} = 3.078 \times 2.4495 \approx 7.540 \text{ ms}.
	\end{align}
	
	\textbf{Parte 3: Comparación e implicaciones técnicas.}
	El umbral de Bonferroni ($7.540$ ms) es \textbf{casi un $50\%$ más severo y amplio} que el umbral de Fisher LSD ($5.046$ ms). 
	\begin{itemize}
		\item \textbf{LSD de Fisher} es una herramienta de alta sensibilidad experimental (bajo error Tipo II): detectará fácilmente diferencias sutiles entre protocolos (como $5.2$ ms), pero al no ajustar por las 10 comparaciones de la familia, su probabilidad real de declarar al menos un falso positivo en todo el experimento asciende a más del $35\%$.
		\item \textbf{Bonferroni} garantiza una especificidad férrea ($\text{FWER} \leq 5\%$), impidiendo falsas alarmas. Sin embargo, su rigidez cobra un precio en potencia estadística (elevado error Tipo II): diferencias reales y operativamente valiosas entre $5.1$ ms y $7.5$ ms pasarán inadvertidas como no significativas. En analítica computacional, cuando $k \geq 5$, se prefiere el término medio equilibrado que ofrece HSD de Tukey.
	\end{itemize}
\end{solucion}

\begin{solucion}
	\begin{enumerate}
		\item \textbf{Cálculo de la gran media general ponderada ($\bar{Y}_{..}$):}
		En diseños con tamaños de muestra heterogéneos ($n_j \neq n$), el promedio general no es simple, sino el promedio ponderado por el número de réplicas grupales ($N = \sum n_j = 6 + 8 + 5 = 19$):
		\begin{align}
			\bar{Y}_{..} = \frac{\sum_{j=1}^3 n_j \bar{Y}_{j.}}{N} = \frac{6(45.0) + 8(52.0) + 5(38.0)}{19} = \frac{270.0 + 416.0 + 190.0}{19} = \frac{876.0}{19} \approx 46.1053\%.
		\end{align}
		
		\item \textbf{Cálculo de sumas de cuadrados, Cuadrados Medios y prueba ANOVA:}
		\textbf{SCTR (desbalanceada):}
		\begin{align}
			\text{SCTR} &= \sum_{j=1}^3 n_j (\bar{Y}_{j.} - \bar{Y}_{..})^2 = 6(45.0 - 46.1053)^2 + 8(52.0 - 46.1053)^2 + 5(38.0 - 46.1053)^2 \nonumber \\
			&= 6(-1.1053)^2 + 8(5.8947)^2 + 5(-8.1053)^2 = 6(1.2217) + 8(34.7475) + 5(65.6959) \nonumber \\
			&= 7.3302 + 277.9800 + 328.4795 = 613.7897 \text{ (\%}^2\text{).}
		\end{align}
		
		\textbf{SCE (desbalanceada):} Suma ponderada de varianzas individuales:
		\begin{align}
			\text{SCE} = \sum_{j=1}^3 (n_j - 1)S_j^2 = (6-1)(12.0) + (8-1)(15.0) + (5-1)(10.0) = 5(12) + 7(15) + 4(10) = 60 + 105 + 40 = 205.000 \text{ (\%}^2\text{).}
		\end{align}
		
		La Suma Total es $\text{SCT} = 613.7897 + 205.000 = 818.7897$.
		
		Grados de libertad: $\nu_{\text{TR}} = 3-1 = 2$, $\nu_E = 19 - 3 = 16$.
		Cuadrados Medios:
		\begin{align}
			\text{CMTR} &= \frac{613.7897}{2} \approx 306.8949, \\
			\text{CME} &= \frac{205.000}{16} = 12.8125.
		\end{align}
		El estadístico experimental resulta:
		\begin{align}
			F = \frac{\text{CMTR}}{\text{CME}} = \frac{306.8949}{12.8125} \approx 23.9528.
		\end{align}
		Con $\alpha = 0.05$ en $\nu_1=2, \nu_2=16$, el umbral es $F_{0.05, \, 2, \, 16} = 3.634$.
		Como $F = 23.953 \gg 3.634$, \textbf{rechazamos $H_0$}. El consumo medio de CPU difiere radicalmente de un entorno de virtualización a otro.
		
		\item \textbf{Cálculo de los coeficientes $R^2$ y $R_{\text{ajustado}}^2$:}
		El \textbf{Coeficiente de Determinación ($R^2$ o Eta-cuadrada $\eta^2$)} mide la fracción total de variabilidad explicada por el factor principal:
		\begin{align}
			R^2 = \frac{\text{SCTR}}{\text{SCT}} = \frac{613.7897}{818.7897} \approx 0.7496 \quad (74.96\%).
		\end{align}
		Esto significa que el $74.96\%$ de toda la fluctuación observada en el consumo de CPU de la granja de servidores es directamente atribuible y explicada por el tipo de entorno de virtualización elegido.
		
		El \textbf{$R^2$ ajustado} penaliza el número de grados de libertad consumidos por los tratamientos respecto al total:
		\begin{align}
			R_{\text{ajustado}}^2 = 1 - \left[ \frac{\text{CME}}{\text{CMT}} \right] = 1 - \left[ \frac{\text{SCE}/\nu_E}{\text{SCT}/\nu_T} \right] = 1 - \left[ \frac{12.8125}{818.7897 / 18} \right] = 1 - \left[ \frac{12.8125}{45.4883} \right] = 1 - 0.28166 \approx 0.7183 \quad (71.83\%).
		\end{align}
	\end{enumerate}
\end{solucion}

\subsubsection*{3. Nivel Analítico (Avanzados / De Conexión)}

\begin{solucion}
	La afirmación del analista es \textbf{peligrosamente incompleta en el escenario analizado}. Es un hecho matemático que el ANOVA goza de robustez estructural frente a desviaciones leves a moderadas de homoscedasticidad cuando las muestras están balanceadas ($n_1 = n_2 = n_3$). Sin embargo, dicha robustez \textbf{no soporta disparidades extremas} en las varianzas grupales.
	
	En nuestro caso, la varianza del Tratamiento 3 ($S_3^2 = 8.9$) es más de \textbf{siete veces superior} a la varianza del Tratamiento 1 ($S_1^2 = 1.2$, cociente $8.9 / 1.2 = 7.41$). Esta heterogeneidad severa indica que el Tratamiento 3 opera con una inestabilidad sistémica muy superior o contiene valores atípicos severos.
	
	\textbf{Procedimiento de verificación y acción metodológica rigurosa:}
	\begin{enumerate}
		\item \textbf{Auditoría cualitativa y cuantitativa de residuos:} Se debe construir un gráfico de residuos estudiantizados e inspeccionar si la dispersión adopta forma de embudo o trompeta (indicio clásico de varianza no constante dependiente de la media).
		\item \textbf{Prueba formal de Levene o Brown-Forsythe:} Se debe someter la muestra a la prueba de Levene ($\alpha = 0.05$). Ante una disparidad de razón $7.4:1$, la prueba de Levene rechazará de forma contundente la hipótesis nula $H_0: \sigma_1^2 = \sigma_2^2 = \sigma_3^2$.
		\item \textbf{Alternativa analítica adecuada:} Dado el fallo confirmado de homoscedasticidad, el analista debe desestimar el estadístico $F$ convencional y emplear el \textbf{ANOVA de Welch (o prueba $F$ de Welch)}, que pondera inversamente los cuadrados medios grupales por las varianzas individuales $S_i^2/n_i$ y ajusta los grados de libertad del denominador, o bien aplicar una transformación logarítmica ($\log Y_{ij}$) si los datos son positivos y asimétricos a la derecha.
	\end{enumerate}
\end{solucion}

\begin{solucion}
	\textbf{Paso 1: Estimación del tamaño del efecto bajo la configuración más adversa.}
	Para un número fijo de tratamientos $k = 4$ con una diferencia máxima especificada $\Delta = 5$ ms y una dispersión de error $\sigma = 4$ ms, la configuración poblacional más difícil de detectar (la que minimiza la suma de cuadrados de los tratamientos $\sum \tau_i^2$) ocurre cuando dos tratamientos están en los extremos distanciados por $\Delta$ y los otros $k-2$ tratamientos se ubican exactamente en el punto medio de ambos.
	
	Bajo esta configuración mínima, la varianza poblacional de los tratamientos es:
	\begin{align}
		\sigma_{\tau}^2 = \frac{\Delta^2}{2k} = \frac{(5)^2}{2(4)} = \frac{25}{8} = 3.125 \text{ ms}^2.
	\end{align}
	
	\textbf{Paso 2: Cálculo del tamaño del efecto estandarizado $f$ de Cohen.}
	El efecto global de Cohen para ANOVA se define como:
	\begin{align}
		f = \frac{\sigma_{\tau}}{\sigma} = \sqrt{\frac{3.125}{16}} = \sqrt{0.1953} \approx 0.4419.
	\end{align}
	Un valor $f = 0.44$ representa un tamaño del efecto de magnitud \textbf{grande a mediana} en la escala estandarizada de Cohen ($f=0.10$ pequeño, $0.25$ mediano, $0.40$ grande).
	
	\textbf{Paso 3: Cálculo del parámetro de no centralidad del ANOVA ($\lambda$) o uso de aproximación de muestras por grupo.}
	La relación teórica entre el número de réplicas por grupo $n$, el efecto de Cohen $f$, y el parámetro de no centralidad del estadístico $F$ es $\lambda = n \cdot k \cdot f^2$.
	
	Para alcanzar un poder estadístico del $80\%$ ($1 - \beta = 0.80$) con $\alpha = 0.05$ y $\nu_1 = k - 1 = 3$ grados de libertad en el numerador, las tablas de potencia computacional (o el paquete \texttt{statsmodels.stats.power.FTestAnovaPower} en Python) dictan que el parámetro de no centralidad requerido debe rondar el valor $\lambda \approx 10.90$.
	
	Despejando $n$ de la relación fundamental de no centralidad:
	\begin{align}
		n = \frac{\lambda}{k \cdot f^2} = \frac{10.90}{4 \times (0.4419)^2} = \frac{10.90}{4 \times 0.1953} = \frac{10.90}{0.7812} \approx 13.95 \text{ réplicas por lenguaje}.
	\end{align}
	
	\textbf{Conclusión de diseño experimental:} Redondeando al entero superior inmediato para garantizar el cumplimiento estricto del poder deseado, se requiere asignar un tamaño de muestra de \textbf{$n = 14$ réplicas por cada lenguaje de backend}, lo que implica un tamaño total de experimento de \textbf{$N = k \times n = 4 \times 14 = 56$ peticiones de API auditadas}.
\end{solucion}

\subsubsection*{4. Nivel Desafiante (Expertos / Deducción Profunda)}

\begin{solucion}
	\begin{enumerate}
		\item \textbf{Demostración algebraica de la partición cuadrática e identidad del doble producto:}
		Sea la identidad de descomposición observacional exacta:
		\begin{align}
			Y_{ij} - \bar{Y}_{..} = (\bar{Y}_{i.} - \bar{Y}_{..}) + (Y_{ij} - \bar{Y}_{i.}).
		\end{align}
		Elevamos ambos miembros al cuadrado:
		\begin{align}
			(Y_{ij} - \bar{Y}_{..})^2 = (\bar{Y}_{i.} - \bar{Y}_{..})^2 + (Y_{ij} - \bar{Y}_{i.})^2 + 2(\bar{Y}_{i.} - \bar{Y}_{..})(Y_{ij} - \bar{Y}_{i.}).
		\end{align}
		Aplicamos el operador suma doble sobre el índice de tratamientos ($i=1,\dots,k$) y el índice de réplicas ($j=1,\dots,n$):
		\begin{align}
			\sum_{i=1}^k \sum_{j=1}^n (Y_{ij} - \bar{Y}_{..})^2 = \sum_{i=1}^k \sum_{j=1}^n (\bar{Y}_{i.} - \bar{Y}_{..})^2 + \sum_{i=1}^k \sum_{j=1}^n (Y_{ij} - \bar{Y}_{i.})^2 + 2\sum_{i=1}^k \sum_{j=1}^n (\bar{Y}_{i.} - \bar{Y}_{..})(Y_{ij} - \bar{Y}_{i.}).
		\end{align}
		En el primer sumando de la derecha, el término $(\bar{Y}_{i.} - \bar{Y}_{..})^2$ es una constante respecto a la suma interior en $j$, por lo que $\sum_{j=1}^n (\bar{Y}_{i.} - \bar{Y}_{..})^2 = n(\bar{Y}_{i.} - \bar{Y}_{..})^2$. Esto constituye por definición la Suma de Cuadrados de Tratamientos ($\text{SCTR}$). El segundo sumando es la Suma de Cuadrados del Error residual ($\text{SCE}$).
		
		Analicemos el término del doble producto cruzado:
		\begin{align}
			T_{\text{cruzado}} = \sum_{i=1}^k (\bar{Y}_{i.} - \bar{Y}_{..}) \left[ \sum_{j=1}^n (Y_{ij} - \bar{Y}_{i.}) \right].
		\end{align}
		Evaluando la suma interior en los corchetes para un tratamiento fijo $i$:
		\begin{align}
			\sum_{j=1}^n (Y_{ij} - \bar{Y}_{i.}) = \sum_{j=1}^n Y_{ij} - \sum_{j=1}^n \bar{Y}_{i.} = n\bar{Y}_{i.} - n\bar{Y}_{i.} = 0.
		\end{align}
		Por una propiedad de la media aritmética, la suma de las desviaciones de un conjunto de datos respecto a su media grupal es estrictamente idéntica a cero. En consecuencia, al multiplicar por el factor constante en $i$, toda la suma doble cruzada se anula algebraicamente sin residuo:
		\begin{align}
			2 \sum_{i=1}^k (\bar{Y}_{i.} - \bar{Y}_{..}) [0] \equiv 0 \implies \text{SCT} = \text{SCTR} + \text{SCE}.
		\end{align}
		
		\item \textbf{Distribuciones Chi-cuadrada bajo la hipótesis nula universal ($H_0: \tau_i = 0$):}
		Bajo $H_0$, todas las variables aleatorias observadas son idénticamente distribuidas $Y_{ij} \sim N(\mu, \sigma^2)$.
		Para el error residual, en cada tratamiento individual $i$, por el Teorema de Gosset sobre la varianza muestral de observaciones normales independientes:
		\begin{align}
			\frac{\sum_{j=1}^n (Y_{ij} - \bar{Y}_{i.})^2}{\sigma^2} = \frac{(n-1)S_i^2}{\sigma^2} \sim \chi^2_{n-1}.
		\end{align}
		Como las $k$ muestras grupales son totalmente independientes entre sí, la suma de variables Chi-cuadrada independientes sigue otra distribución Chi-cuadrada cuyos grados de libertad son la suma directa de los sumandos:
		\begin{align}
			\frac{\text{SCE}}{\sigma^2} = \sum_{i=1}^k \left[ \frac{(n-1)S_i^2}{\sigma^2} \right] \sim \chi^2_{\sum_{i=1}^k (n-1)} = \chi^2_{k(n-1)} = \chi^2_{N-k}.
		\end{align}
		Asimismo, bajo $H_0$, las $k$ medias muestrales son independientes y se distribuyen $\bar{Y}_{i.} \sim N\left(\mu, \frac{\sigma^2}{n}\right)$. Aplicando de nuevo el Teorema de Gosset a este vector de medias grupales de tamaño $k$ con media general $\bar{Y}_{..}$:
		\begin{align}
			\frac{\sum_{i=1}^k (\bar{Y}_{i.} - \bar{Y}_{..})^2}{\sigma^2 / n} = \frac{n \sum_{i=1}^k (\bar{Y}_{i.} - \bar{Y}_{..})^2}{\sigma^2} = \frac{\text{SCTR}}{\sigma^2} \sim \chi^2_{k-1}.
		\end{align}
		
		\item \textbf{Invocación del Teorema de Cochran y deducción del estadístico $F$:}
		El \textbf{Teorema de William G. Cochran} establece formalmente en álgebra lineal que: si un vector normal multivariado estándar $\mathbf{Z} \sim N(\mathbf{0}, \mathbf{I}_N)$ satisface la partición en formas cuadráticas $\mathbf{Z}^T \mathbf{Z} = \sum_{m=1}^M \mathbf{Z}^T \mathbf{A}_m \mathbf{Z}$, donde las matrices $\mathbf{A}_m$ son simétricas con rangos $r_m = \text{rango}(\mathbf{A}_m)$, entonces \textbf{las formas cuadráticas $\mathbf{Z}^T \mathbf{A}_m \mathbf{Z}$ son variables aleatorias Chi-cuadrada mutuamente e incondicionalmente independientes con $r_m$ grados de libertad, si y solo si la suma de los rangos iguala la dimensión total del espacio: $\sum_{m=1}^M r_m = N$}.
		
		En nuestro diseño unifactorial, el vector de desviaciones estandarizadas satisface la partición exacta en matrices simétricas de proyección:
		\begin{align}
			\frac{\text{SCT}}{\sigma^2} = \frac{\text{SCTR}}{\sigma^2} + \frac{\text{SCE}}{\sigma^2}.
		\end{align}
		La dimensión total del espacio de desviaciones respecto a la gran media es $\nu_T = N - 1$. Los rangos del lado derecho son $\nu_{\text{TR}} = k - 1$ para el tratamiento y $\nu_E = N - k$ para el error. Sumando las dimensiones de los dos subespacios proyectivos:
		\begin{align}
			(k - 1) + (N - k) = N - 1 = \nu_T.
		\end{align}
		Al cumplirse con rigurosa exactitud la condición de partición dimensional del Teorema de Cochran, las matrices proyectoras de $\text{SCTR}$ y $\text{SCE}$ son \textbf{idempotentes y mutuamente ortogonales ($\mathbf{A}_{\text{TR}} \mathbf{A}_E = \mathbf{0}$)}. Por tanto, las variables $\text{CMTR}$ y $\text{CME}$ son estocástica y funcionalmente independientes.
		
		En consecuencia, el cociente normalizado de dos independientes chi-cuadrada divididas cada una entre sus respectivos grados de libertad define, por la propia axiomática del cálculo de probabilidades, a la distribución $F$ de Snedecor central:
		\begin{align}
			F = \frac{\dfrac{\text{SCTR}/\sigma^2}{k-1}}{\dfrac{\text{SCE}/\sigma^2}{N-k}} = \frac{\text{CMTR}}{\text{CME}} \sim F_{k-1, \, N-k}.
		\end{align}
	\end{enumerate}
\end{solucion}

\begin{solucion}
	\begin{enumerate}
		\item \textbf{Deducción sistemática del valor esperado de los Cuadrados Medios ($E[\text{CM}]$):}
		Sea el modelo en bloques completos al azar $Y_{ij} = \mu + \tau_i + \beta_j + \varepsilon_{ij}$.
		La esperanza de cualquier observación individual es $E[Y_{ij}] = \mu + \tau_i + \beta_j$, ya que $E[\varepsilon_{ij}] = 0$.
		La media grupal del tratamiento $i$ sobre los $b$ bloques es:
		\begin{align}
			\bar{Y}_{i.} = \frac{1}{b}\sum_{j=1}^b (\mu + \tau_i + \beta_j + \varepsilon_{ij}) = \mu + \tau_i + \frac{1}{b}\sum_{j=1}^b \beta_j + \bar{\varepsilon}_{i.} = \mu + \tau_i + 0 + \bar{\varepsilon}_{i.} = \mu + \tau_i + \bar{\varepsilon}_{i.},
		\end{align}
		y de forma simétrica para el bloque $j$ sobre los $k$ tratamientos, $\bar{Y}_{.j} = \mu + \beta_j + \bar{\varepsilon}_{.j}$, mientras que la gran media es $\bar{Y}_{..} = \mu + \bar{\varepsilon}_{..}$.
		
		Analicemos el Cuadrado Medio del Tratamiento ($\text{CMTR} = \dfrac{b}{k-1}\sum_{i=1}^k (\bar{Y}_{i.} - \bar{Y}_{..})^2$):
		Sustituyendo las expresiones reducidas de las medias:
		\begin{align}
			\bar{Y}_{i.} - \bar{Y}_{..} = (\mu + \tau_i + \bar{\varepsilon}_{i.}) - (\mu + \bar{\varepsilon}_{..}) = \tau_i + (\bar{\varepsilon}_{i.} - \bar{\varepsilon}_{..}).
		\end{align}
		Elevando al cuadrado e invocando la esperanza sobre la suma:
		\begin{align}
			E\left[ \sum_{i=1}^k (\bar{Y}_{i.} - \bar{Y}_{..})^2 \right] = E\left[ \sum_{i=1}^k \tau_i^2 + 2\sum_{i=1}^k \tau_i (\bar{\varepsilon}_{i.} - \bar{\varepsilon}_{..}) + \sum_{i=1}^k (\bar{\varepsilon}_{i.} - \bar{\varepsilon}_{..})^2 \right].
		\end{align}
		Como $\tau_i$ es un parámetro fijo no aleatorio independiente del error, y $E[\bar{\varepsilon}_{i.} - \bar{\varepsilon}_{..}] = 0$, el término del medio se anula. Para el tercer término, por ser un promedio aleatorio de $b$ observaciones independientes con varianza $\sigma^2/b$:
		\begin{align}
			E\left[ \sum_{i=1}^k (\bar{\varepsilon}_{i.} - \bar{\varepsilon}_{..})^2 \right] = (k - 1) \text{Var}(\bar{\varepsilon}_{i.}) = (k - 1) \frac{\sigma^2}{b}.
		\end{align}
		Sustituyendo y multiplicando por la constante externa $\dfrac{b}{k-1}$ de $\text{CMTR}$:
		\begin{align}
			E[\text{CMTR}] = \frac{b}{k-1} \left[ \sum_{i=1}^k \tau_i^2 + \frac{(k-1)\sigma^2}{b} \right] = \sigma^2 + \frac{b}{k-1}\sum_{i=1}^k \tau_i^2.
		\end{align}
		Por el principio de simetría exacta sobre los índices y las ligaduras de los bloques $\sum \beta_j = 0$, la derivación para el Cuadrado Medio del Bloque es estructuralmente gemela:
		\begin{align}
			E[\text{CMB}] = \sigma^2 + \frac{k}{b-1}\sum_{j=1}^b \beta_j^2.
		\end{align}
		Para el Cuadrado Medio del Error ($\text{CME}$), el residuo en la celda $(i,j)$ del DBCA es:
		\begin{align}
			e_{ij} = Y_{ij} - \bar{Y}_{i.} - \bar{Y}_{.j} + \bar{Y}_{..} = \varepsilon_{ij} - \bar{\varepsilon}_{i.} - \bar{\varepsilon}_{.j} + \bar{\varepsilon}_{..}.
		\end{align}
		Al elevar al cuadrado y sumar sobre las $kb$ celdas, los parámetros deterministas $\tau_i$ y $\beta_j$ desaparecen en su totalidad por cancelación cruzada. El operador esperanza sobre esta suma puramente aleatoria de dimensión $(k-1)(b-1)$ arroja:
		\begin{align}
			E[\text{SCE}] = (k - 1)(b - 1) \sigma^2 \implies E[\text{CME}] = \frac{E[\text{SCE}]}{(k-1)(b-1)} = \sigma^2.
		\end{align}
		El $\text{CME}$ es el único estimador perfectamente insesgado de la varianza intrínseca del error $\sigma^2$.
		
		\item \textbf{Deducción de la Eficiencia Relativa ($\text{ER}$) del DBCA respecto al DCA:}
		Si el diseño no se hubiera bloqueado y los datos se hubieran procesado como un DCA de un solo factor con $k$ tratamientos de tamaño $b$ ($N=kb$), la Suma de Cuadrados del Error en el DCA habría sido la unión del residuo intrínseco con la variabilidad que el DBCA separó como bloques:
		\begin{align}
			\text{SCE}_{\text{DCA}} = \text{SCB} + \text{SCE}_{\text{DBCA}}.
		\end{align}
		Los grados de libertad del error del DCA no bloqueado habrían sido $\nu_{E,\text{DCA}} = N - k = k(b - 1)$.
		Evaluando la esperanza o estimador de la varianza del error en ese DCA hipotético ($S_{\text{DCA}}^2 = \dfrac{\text{SCE}_{\text{DCA}}}{k(b-1)}$):
		\begin{align}
			S_{\text{DCA}}^2 = \frac{\text{SCB} + \text{SCE}_{\text{DBCA}}}{k(b - 1)} = \frac{(b-1)\text{CMB} + (k-1)(b-1)\text{CME}}{k(b-1)}.
		\end{align}
		Para ajustar esta aproximación por la diferencia finita de grados de libertad al reasignar el diseño en muestras pequeñas (el método exacto de Kempthorne y Fisher para eficiencia en diseños aleatorizados), se distribuyen los grados de libertad totales del error e intragrupales en la combinación convexa ponderada por grados de libertad exactos:
		\begin{align}
			S_{\text{DCA}}^2 \approx \frac{(b-1)\text{CMB} + b(k-1)\text{CME}}{(b-1) + b(k-1)} = \frac{(b-1)\text{CMB} + b(k-1)\text{CME}}{kb - 1}.
		\end{align}
		La \textbf{Eficiencia Relativa ($\text{ER}$)} se define como la razón de las varianzas experimentales del error ($S_{\text{DCA}}^2 / S_{\text{DBCA}}^2 = S_{\text{DCA}}^2 / \text{CME}$):
		\begin{align}
			\text{ER} = \frac{S_{\text{DCA}}^2}{\text{CME}} = \frac{(b-1)\dfrac{\text{CMB}}{\text{CME}} + b(k-1)}{kb - 1} = \frac{(b-1)F_{\text{Bloque}} + b(k-1)}{kb - 1}.
		\end{align}
		
		\textbf{Condición analítica de superioridad del bloqueo ($\text{ER} > 1$):}
		Para que el Diseño en Bloques supere en precisión al diseño unifactorial aleatorizado, se requiere $\text{ER} > 1$:
		\begin{align}
			\frac{(b-1)\text{CMB} + b(k-1)\text{CME}}{kb - 1} > \text{CME} \implies (b-1)\text{CMB} + b(k-1)\text{CME} > (kb - 1)\text{CME}.
		\end{align}
		Expandiendo el producto algebraico en el lado derecho:
		\begin{align}
			(b-1)\text{CMB} + (kb - b)\text{CME} > (kb - 1)\text{CME} \implies (b-1)\text{CMB} > (kb - 1 - kb + b)\text{CME} = (b-1)\text{CME}.
		\end{align}
		Dividiendo ambos lados entre la cantidad estrictamente positiva $(b-1)$:
		\begin{align}
			\text{CMB} > \text{CME} \iff F_{\text{Bloque}} = \frac{\text{CMB}}{\text{CME}} > 1.
		\end{align}
		\textbf{Conclusión teórica:} El Diseño en Bloques Completos al Azar \textbf{aporta una ganancia estricta y real de eficiencia en estimación ($\text{ER} > 1$) siempre que el Cuadrado Medio del Bloque supere al Cuadrado Medio del Error ($F_{\text{Bloque}} > 1$)}. Si el factor de bloqueo fue irrelevante o mal escogido de modo que $\text{CMB} < \text{CME}$, el bloqueo habrá sido un lastre analítico que habrá destruido innecesariamente grados de libertad del error sin absorber varianza de la infraestructura.
	\end{enumerate}
\end{solucion}
